\documentclass[12pt]{article}
\usepackage[utf8]{inputenc}
\usepackage[T1]{fontenc}
\usepackage{hyperref}
\title{Tunebo, Angosturas Language Resources for K-12 Curriculum Development}
\author{Compiled by Claude Sonnet 5 (claude-sonnet-5) for Andr\'as Kornai}
\date{}
\begin{document}
\maketitle

\section{Introduction}
Angosturas Tunebo (ISO 639-3 tnd; Glottocode ango1257; also called Bahiyakuwa) is a Chibchan language of Arauca department, Colombia, near the headwaters of the Casanare River. It is one of several Tunebo (U'wa) varieties. A search summary gave about 50 speakers according to Ethnologue 2016. I did not check Ethnologue directly. No material specific to this variety appeared in the searches. What exists is general Tunebo/U'wa documentation.

\section{Primary Educational Resources}
No primers, textbooks or school programmes for Angosturas Tunebo were found.

\section{Linguistic and Reference Materials}
\begin{itemize}
\item Wikipedia, Uwa language (covers the Tunebo varieties as a group): \url{https://en.wikipedia.org/wiki/Uwa_language}
\item Chibchan documentation PDF on Tunebo, University of Missouri (title only seen in search results; whether it covers Angosturas is unverified): \url{https://dice.missouri.edu/assets/docs/chibchan/Tunebo.pdf}
\item Grambank entry for ango1257: \url{https://grambank.clld.org/languages/ango1257}
\item WALS entry for Tunebo (tnb): \url{https://wals.info/languoid/lect/wals_code_tnb}
\item Encyclopedia.com entry on the Tunebo people: \url{https://www.encyclopedia.com/humanities/encyclopedias-almanacs-transcripts-and-maps/tunebo}
\end{itemize}

\section{Religious and Community Resources}
\begin{itemize}
\item Joshua Project, Western Tunebo in Colombia (a different variety, not Angosturas): \url{https://joshuaproject.net/people_groups/15387}
\end{itemize}

\section{Digital and Audio Resources}
\begin{itemize}
\item Endangered Languages Project page for the Tunebo/U'wa language (rated threatened): \url{https://endangeredlanguages.com/lang/3527}
\end{itemize}
I found no audio or corpus specific to Angosturas.

\section{Curriculum Development Context}
Material for this variety is very thin. Tunebo varieties differ enough that resources for the others (Western, Central, Barro Negro) may not transfer directly. Any curriculum would need documentation work with the remaining speakers, who are mainly older people, and coordination with U'wa education efforts.
\end{document}